%============================================================
% CH2 · PATTERN 1 — CYCLE DETECTION AND RECOVERY: THEORY
%============================================================
\subsection{Pattern Theory}
\label{ssec:ch2-p1-theory}

\begin{patternbox}{Pattern 1 at a Glance}
Run DFS from every white vertex. Keep the vertices of the current path \textbf{gray}. The first time DFS, sitting at $u$, examines an edge to a vertex $v$ that is \textbf{still on the path} (gray; in an undirected graph: and not through the edge it came in by), the path from $v$ down to $u$ plus that edge is a cycle. If the search ends without such an edge, there is no cycle.

\begin{center}
\begin{tikzpicture}[x=1.25cm]
    \foreach \v/\i/\s in {s/0/gvsrc, a/1/gvseen, v/2/gvsel, b/3/gvsel, c/4/gvsel, u/5/gvcur}
        \node[\s] (p\i) at (\i,0) {$\v$};
    \foreach \i [evaluate=\i as \j using int(\i+1)] in {0,...,4} \draw[gtree, -{Stealth[length=2.4mm]}] (p\i) -- (p\j);
    \draw[gback] (p5.north) to[bend right=40] node[lbl, above]{back edge $u \to v$} (p2.north);
    \node[gvdone] (x) at (1.5,-1.3) {$x$};
    \draw[gtree, -{Stealth[length=2.4mm]}] (p1) -- (x);
    \draw[gcross] (p4) to[bend left=12] node[lbl, below, pos=0.6]{to black: harmless} (x);
    \draw[decorate, decoration={brace, amplitude=5pt}, darkgray] (-0.35,0.65) -- (5.35,0.65)
        node[midway, above=5pt, font=\scriptsize]{the stack (gray vertices) $=$ the current path};
    \node[note, text width=3.6cm, align=left, anchor=west] at (6.0,0) {cycle $= v \to b \to c \to u \to v$\\(purple + orange), read off \texttt{parent[]} from $u$ up to $v$};
\end{tikzpicture}
\end{center}
\begin{recurrencebox}
\centering\small
colour $0 \to 1$ when entered, $1 \to 2$ when left;\quad examine $u \to v$:\\[3pt]
$\text{colour}[v] = 0$: tree edge, \texttt{parent}$[v] = u$ \qquad $\text{colour}[v] = 1$: \textbf{cycle} $v \leadsto u \to v$ \qquad $\text{colour}[v] = 2$: ignore\\[3pt]
undirected: skip the edge id \texttt{parentEdge}$[u]$ \qquad total \VE, cycle printed in \bigO{\text{its length}}
\end{recurrencebox}
\end{patternbox}

%------------------------------------------------------------
\subsubsection*{From the lessons to the pattern}

Lesson~\ref{ssec:ch2-cycles} proved that back edges and cycles are the same thing: a graph has a cycle iff a DFS meets a back edge, and the cycle is sitting on the call stack at that very moment. So the pattern is not ``find a clever algorithm'', it is ``run the right DFS and read the answer off the stack''. The difficulty in contest problems lies elsewhere:
\begin{itemize}[leftmargin=*]
    \item recognising that the question \emph{is} about a cycle --- ``a round trip'', ``a deadlock'', ``contradictory orders'', ``an infinite loop'', ``a process that never ends'';
    \item getting the details right for the kind of graph: directed or undirected, simple or with parallel edges and self-loops, connected or not;
    \item printing the cycle in the format the judge wants (often with the first vertex repeated at the end);
    \item and not overflowing the stack on $10^5$ vertices.
\end{itemize}

%------------------------------------------------------------
\subsubsection*{When to recognise it (and when not)}

\begin{itemize}[leftmargin=*]
    \item ``\textbf{find} a cycle / a round trip / a route that returns to its start'' --- print any one;
    \item ``is the order consistent?'', ``can all courses be finished?'', ``do the screenshots contradict each other?'' --- a directed cycle means \emph{no};
    \item ``is the graph a tree / a forest / a cactus / a single cycle?'' --- cycle tests per component (often by counting, $m$ vs.\ $n - c$);
    \item ``does the simulation run forever?'' when the next step is \emph{not} forced (if it is forced, it is a functional graph: Chapter~1, Pattern~5).
\end{itemize}
\textbf{Not this pattern:} the \emph{shortest} cycle needs BFS from every vertex (Chapter~3, and Chapter~22 for faster methods); an \emph{odd} cycle is a bipartiteness test (Chapter~4); a cycle of negative weight is Bellman--Ford (Chapter~8); a cycle through \emph{every} vertex or \emph{every} edge is Hamiltonian (Chapter~23) or Eulerian (Chapter~21); counting all cycles is exponential in general.

%------------------------------------------------------------
\subsubsection*{The algorithm and its invariant}

\begin{invariantbox}[Why reading the stack gives a cycle]
\textbf{Invariant.} At every moment, the gray vertices are exactly the vertices on the stack, in order $x_0, x_1, \dots, x_k$ with $\texttt{parent}[x_{i+1}] = x_i$, and every edge $x_i \to x_{i+1}$ is a real edge of the graph.

\emph{Detection is sound.} If the edge $u \to v$ with $u = x_k$ finds $v$ gray, then $v = x_j$ for some $j \le k$, and $x_j \to x_{j+1} \to \dots \to x_k \to x_j$ is a closed walk without repeated vertices --- a cycle. In an undirected graph, $j = k-1$ is impossible unless the edge is a \emph{second} edge between $x_{k-1}$ and $x_k$, because the tree edge itself is skipped; so in a simple graph the cycle has at least 3 vertices.

\emph{Detection is complete.} By the cycle theorem of Lesson~\ref{ssec:ch2-cycles}, if any cycle exists, the DFS meets a back edge before it finishes.

\emph{Cost.} Every vertex is pushed once and every adjacency entry examined once: \VE. The parent walk visits each cycle vertex once.
\end{invariantbox}

\begin{diagrambox}
\begin{tikzpicture}[x=1cm, y=1cm]
    % directed example
    \begin{scope}
        \node[font=\small\bfseries] at (1.6,2.2) {directed: gray or black?};
        \node[gvseen] (a) at (0,0.6) {1};
        \node[gvseen] (b) at (1.6,0.6) {2};
        \node[gvcur] (c) at (3.2,0.6) {3};
        \node[gvdone] (d) at (1.6,-0.8) {4};
        \draw[gtree, -{Stealth[length=2.4mm]}] (a) -- (b);
        \draw[gtree, -{Stealth[length=2.4mm]}] (b) -- (c);
        \draw[gtree, -{Stealth[length=2.4mm]}] (a) -- (d);
        \draw[gcross] (c) -- node[lbl, right, pos=0.4]{black} (d);
        \draw[gback] (c) to[bend right=50] node[lbl, above]{gray} (a);
        \node[font=\scriptsize, text width=4.4cm, align=center] at (1.6,-1.85) {$3 \to 4$ meets a finished vertex: no cycle through it. $3 \to 1$ meets the stack: cycle $1, 2, 3$.};
    \end{scope}
    % undirected example
    \begin{scope}[xshift=6.4cm]
        \node[font=\small\bfseries] at (1.6,2.2) {undirected: which edge to skip?};
        \node[gvseen] (p) at (0,0.6) {1};
        \node[gvseen] (q) at (1.6,0.6) {2};
        \node[gvcur] (r) at (3.2,0.6) {3};
        \draw[gtree] (p) -- node[gw, above]{e0} (q);
        \draw[gtree] (q) -- node[gw, above]{e1} (r);
        \draw[gfaded, dashed] (r) to[bend left=35] node[gw, below]{e1 again} (q);
        \draw[gback, -] (r) to[bend right=55] node[lbl, above]{e2} (p);
        \node[font=\scriptsize, text width=4.6cm, align=center] at (1.6,-1.85) {At 3, the edge e1 (grey) leads back to the parent: skipped. e2 leads to the gray 1: cycle $1, 2, 3$.};
    \end{scope}
\end{tikzpicture}
\figcap{The two decisions of the pattern. Directed: only gray targets close cycles. Undirected: skip exactly the edge used to enter the current vertex.}
\end{diagrambox}

%------------------------------------------------------------
\subsubsection*{Printing the cycle}

Judges usually want the cycle as a list of vertices in walking order, often with the start repeated at the end (``$k$, then $k$ cities, first $=$ last''). With parent pointers:

\begin{diagrambox}
\begin{tikzpicture}[x=1.15cm]
    \node[font=\scriptsize\bfseries, anchor=east] at (-0.4,0) {stack};
    \foreach \v/\i/\s in {1/0/gvseen, 4/1/gvseen, 7/2/gvsel, 2/3/gvsel, 9/4/gvsel, 5/5/gvcur}
        \node[\s] (s\i) at (\i,0) {\v};
    \foreach \i [evaluate=\i as \j using int(\i+1)] in {0,...,4} \draw[gtree, -{Stealth[length=2.4mm]}] (s\i) -- (s\j);
    \draw[gback] (s5.north) to[bend right=40] (s2.north);
    \node[font=\scriptsize, anchor=west] at (6.4,0.35) {back edge $5 \to 7$};
    \node[font=\scriptsize\ttfamily, align=left, anchor=west] at (0,-1.1) {walk parents from 5 to 7:\ \ 5 9 2 7};
    \node[font=\scriptsize\ttfamily, align=left, anchor=west] at (0,-1.6) {reverse:\ \ \ \ \ \ \ \ \ \ \ \ \ \ \ \ \ \ \ 7 2 9 5};
    \node[font=\scriptsize\ttfamily, align=left, anchor=west] at (0,-2.1) {append the start again:\ \ \ \ 7 2 9 5 7\ \ \ (k = 5)};
    \node[note, text width=4.6cm, anchor=west, align=left] at (6.4,-1.6) {Undirected graphs need no reversal: the cycle is the same read in either direction. Directed graphs do --- \texttt{5 9 2 7} walks the edges backwards.};
\end{tikzpicture}
\end{diagrambox}

%------------------------------------------------------------
\subsubsection*{Variations map}

\begin{diagrambox}
\begin{tikzpicture}[
    var/.style={rectangle, rounded corners=2pt, draw=patternborder, fill=patternbg,
                font=\scriptsize, align=center, text width=3.25cm, minimum height=1.3cm, inner sep=3pt}]
    \node[var, fill=answerfill, draw=accentorange, text width=5cm, font=\scriptsize\bfseries] (root) at (0,0)
        {DFS with colours; a back edge closes a cycle};
    \node[var] (v1) at (-5.4,-2) {\textbf{undirected, simple}\\skip the parent, $\ge 3$ vertices\\(Round Trip, 2.1.1)};
    \node[var] (v2) at (-1.8,-2) {\textbf{directed}\\gray, not visited\\(Round Trip II, 2.1.2)};
    \node[var] (v3) at (1.8,-2)  {\textbf{one cycle per component?}\\count: $m = n$, connected\\(E1.2)};
    \node[var] (v4) at (5.4,-2)  {\textbf{grid cycles}\\same colour, length $\ge 4$\\(E1.3)};
    \node[var] (v5) at (-5.4,-3.6) {\textbf{which components are acyclic?}\\trees among components\\(E1.4)};
    \node[var] (v6) at (-1.8,-3.6) {\textbf{consistency of orders}\\constraints $\to$ arrows, cycle $=$ no\\(E1.5)};
    \node[var] (v7) at (1.8,-3.6)  {\textbf{which edge to delete?}\\cycle + in-degree 2\\(E1.6)};
    \node[var] (v8) at (5.4,-3.6)  {\textbf{find the cycle, then use it}\\pseudotrees, entry vertex\\(Practice set)};
    \node[draw=patternborder, dashed, rounded corners=4pt, fit=(v1)(v4)(v5)(v8), inner sep=5pt] (grp) {};
    \draw[flow] (root.south) -- (grp.north);
    \node[note, below=2mm of grp] {shortest / odd / negative / Hamiltonian / Euler cycles are different problems: Chapters 3, 4, 8, 23, 21};
\end{tikzpicture}
\end{diagrambox}

\textbf{Twists and how to handle them}
\begin{itemize}[leftmargin=*]
    \item \textbf{Parallel edges in an undirected graph} are a cycle of length 2; skip the parent \emph{edge id}, not the parent vertex (template 1).
    \item \textbf{Self-loops} are cycles of length 1 in both kinds of graph. If the problem forbids them as answers, check the statement: many inputs simply exclude them.
    \item \textbf{Disconnected input:} the cycle may live in any component --- the outer loop over all start vertices is not optional.
    \item \textbf{``Cycle through vertex $x$'':} a directed cycle through $x$ exists iff some successor of $x$ reaches $x$ (one search in the reversed graph from $x$, Pattern~5); an undirected one iff $x$ has two neighbours that are connected without $x$.
    \item \textbf{Only yes/no, undirected:} count, no search order needed: a cycle exists iff $m > n - c$ (template 3, or a DSU).
    \item \textbf{Only yes/no, directed:} Kahn's peeling of in-degree-0 vertices (Chapter~5) leaves vertices behind iff there is a cycle; the leftovers also let you \emph{print} a cycle (Approach~3 of Problem~\ref{ssec:ch2-round-trip-ii}).
\end{itemize}

%------------------------------------------------------------
\subsubsection*{Code templates}

\begin{lstlisting}[style=cpp]
#include <bits/stdc++.h>
using namespace std;

// (1) UNDIRECTED, multi-edges and self-loops allowed: skip the EDGE we came in by.
//     Returns a cycle v0 v1 ... vk-1 (closing edge vk-1 - v0), or {} for a forest. Iterative.
vector<int> undirectedCycle(int n, const vector<pair<int, int>>& edges) {
    vector<vector<pair<int, int>>> adj(n);             // (neighbour, edge id)
    for (int i = 0; i < (int)edges.size(); i++) {
        adj[edges[i].first].push_back({edges[i].second, i});
        if (edges[i].first != edges[i].second) adj[edges[i].second].push_back({edges[i].first, i});
    }
    vector<int> color(n, 0), par(n, -1), parEdge(n, -1);
    vector<size_t> it(n, 0);
    for (int s = 0; s < n; s++) {
        if (color[s]) continue;
        color[s] = 1;
        vector<int> st(1, s);
        while (!st.empty()) {
            int u = st.back();
            if (it[u] == adj[u].size()) { color[u] = 2; st.pop_back(); continue; }
            int v = adj[u][it[u]].first, id = adj[u][it[u]].second;
            it[u]++;
            if (id == parEdge[u]) continue;            // the edge we came in by
            if (color[v] == 0) { color[v] = 1; par[v] = u; parEdge[v] = id; st.push_back(v); }
            else if (color[v] == 1) {                  // back edge to an ancestor (or a self-loop)
                vector<int> cyc;
                for (int x = u; x != v; x = par[x]) cyc.push_back(x);
                cyc.push_back(v);
                reverse(cyc.begin(), cyc.end());
                return cyc;                            // v ... u, closed by the edge u - v
            }
        }
    }
    return {};
}

// (2) DIRECTED: an edge to a GRAY vertex closes a cycle. Returns v0 -> ... -> vk-1 (-> v0) or {}.
vector<int> directedCycle(const vector<vector<int>>& adj) {
    int n = adj.size();
    vector<int> color(n, 0), par(n, -1);
    vector<size_t> it(n, 0);
    for (int s = 0; s < n; s++) {
        if (color[s]) continue;
        color[s] = 1;
        vector<int> st(1, s);
        while (!st.empty()) {
            int u = st.back();
            if (it[u] == adj[u].size()) { color[u] = 2; st.pop_back(); continue; }
            int v = adj[u][it[u]++];
            if (color[v] == 0) { color[v] = 1; par[v] = u; st.push_back(v); }
            else if (color[v] == 1) {
                vector<int> cyc;
                for (int x = u; x != v; x = par[x]) cyc.push_back(x);
                cyc.push_back(v);
                reverse(cyc.begin(), cyc.end());
                return cyc;
            }
        }
    }
    return {};
}

// (3) Yes/no for an undirected graph without a search order: cycle iff m > n - (#components).
bool undirectedHasCycle(int n, const vector<pair<int, int>>& edges) {
    vector<int> p(n);
    iota(p.begin(), p.end(), 0);
    auto f = [&](int x) {                              // iterative find (no deep recursion)
        while (p[x] != x) x = p[x] = p[p[x]];
        return x;
    };
    for (auto& e : edges) {
        int a = f(e.first), b = f(e.second);
        if (a == b) return true;                       // this edge closes a cycle
        p[a] = b;
    }
    return false;
}

int main() {
    // the graph D of Part A, 0-indexed (vertex k is k-1)
    vector<vector<int>> D = {{1, 4, 5}, {2, 3}, {0}, {2}, {3, 5}, {6}, {5}, {6}};
    for (int x : directedCycle(D)) cout << x + 1 << " ";       // 1 2 3
    cout << "\n";
    // G3 of Chapter 1, 0-indexed
    vector<pair<int, int>> G = {{0, 1}, {0, 3}, {1, 2}, {1, 4}, {3, 4}, {4, 5}};
    for (int x : undirectedCycle(7, G)) cout << x + 1 << " ";   // 1 2 5 4
    cout << "\n";
    // a path 1 - 2 - 3 and two parallel edges between 6 and 7
    vector<pair<int, int>> F = {{0, 1}, {1, 2}, {5, 6}, {6, 5}};
    for (int x : undirectedCycle(7, F)) cout << x + 1 << " ";   // 6 7
    cout << "\n" << undirectedHasCycle(7, G) << undirectedHasCycle(3, {{0, 1}, {1, 2}}) << "\n";  // 10
}
\end{lstlisting}

\begin{insightbox}[Which template when?]
\renewcommand{\arraystretch}{1.25}
\begin{tabularx}{\linewidth}{@{}>{\bfseries}l X@{}}
(1) \texttt{undirectedCycle} & print a cycle of an undirected graph; safe with parallel edges and self-loops \\
(2) \texttt{directedCycle} & print a directed cycle, or prove acyclicity (then the finishing order is a topological order, Chapter~5) \\
(3) \texttt{undirectedHasCycle} & yes/no only, or ``the first edge that closes a cycle'' when edges arrive one by one (Chapter~6) \\
\end{tabularx}
\end{insightbox}

%------------------------------------------------------------
\subsubsection*{Common mistakes}

\begin{warningbox}[Mistake 1 --- ``visited'' instead of ``gray'' in a directed graph]
\begin{lstlisting}[style=snippet]
else if (color[v] != 0) { /* cycle! */ }      // WRONG: a black v is not on the path
else if (color[v] == 1) { /* cycle! */ }      // RIGHT
\end{lstlisting}
Every DAG with a vertex of in-degree $\ge 2$ makes the wrong version report a cycle (and usually crash while printing it, Lesson~\ref{ssec:ch2-cycles}).
\end{warningbox}

\begin{warningbox}[Mistake 2 --- forgetting to skip the parent in an undirected graph]
Without the skip, every tree edge $p - v$ is seen again from $v$ and reported as the ``cycle'' $p, v$ of length 2. Round Trip requires at least 3 cities, so this is \wa{} on almost every test.
\end{warningbox}

\begin{warningbox}[Mistake 3 --- starting only from vertex 1]
The cycle can be in a component that does not contain vertex 1. Always loop over all start vertices.
\end{warningbox}

\begin{warningbox}[Mistake 4 --- recursion on $10^5$ vertices]
A path of $10^5$ cities followed by one closing road is a legal input; a recursive DFS goes $10^5$ levels deep. Use the iterative templates, or make sure the judge has a large stack (Lesson~\ref{ssec:ch2-call-stack}).
\end{warningbox}

\begin{warningbox}[Mistake 5 --- printing the directed cycle backwards]
The parent walk from $u$ to $v$ lists the cycle against the direction of the edges. Reverse it before printing (template 2); a judge that checks every consecutive pair rejects the unreversed list.
\end{warningbox}

\begin{checklistbox}[``I'm stuck'' checklist for Pattern 1]
\begin{itemize}[leftmargin=*]
    \item[$\square$] Directed or undirected? Simple, or are parallel edges and self-loops possible?
    \item[$\square$] Does a cycle mean ``yes'' (round trip exists) or ``no'' (orders contradict, deadlock)?
    \item[$\square$] Do I need the cycle itself, or only its existence (then counting or Kahn may be simpler)?
    \item[$\square$] Am I starting a search from \emph{every} white vertex?
    \item[$\square$] Is the minimum length of a valid cycle (3 cities? 4 cells?) guaranteed by my skip rule?
    \item[$\square$] Output format: number of vertices, first vertex repeated, direction of the walk?
\end{itemize}
\end{checklistbox}
